\documentclass[14 pt]{extarticle}

	\usepackage[frenchb]{babel}
	\usepackage[utf8]{inputenc}  
	\usepackage[T1]{fontenc}
	\usepackage{amssymb}
	\usepackage[mathscr]{euscript}
	\usepackage{stmaryrd}
	\usepackage{amsmath}
	\usepackage{tikz}
	\usepackage[all,cmtip]{xy}
	\usepackage{amsthm}
	\usepackage{varioref}
	\usepackage{geometry}
	\geometry{a4paper}
	\usepackage{lmodern}
	\usepackage{hyperref}
	\usepackage{array}
	 \usepackage{fancyhdr}
\renewcommand{\theenumi}{\alph{enumi})}
	\pagestyle{fancy}
	\theoremstyle{plain}
	\fancyfoot[C]{} 
	\fancyhead[L]{Contrôle}
	\fancyhead[R]{12 janvier 2026}\geometry{
 a4paper,
 total={170mm,257mm},
 left=20mm,
 top=20mm,
 }
	
	
	\title{Interrogation chapitre 5}
	\date{}
	\begin{document}

\begin{center}{\Large Contrôle chapitre 5}\\ 
 \end{center}
 
 
  \subsection*{Exercice 1 (6 points)}
  
Sur votre copie, posez et effectuez les calculs suivants : \begin{enumerate}
\item $56,18 + 34, 9$
\item $56,18 - 34, 9$
\item $5 \times 34, 9$
\item $56,18 \times 34, 9$
\end{enumerate}  

\subsection*{Exercice 2 (4 points)}

Le point culminant de Notre-Dame de Paris est à 96 $m$. Celui de la cathédrale de Strasbourg est de 45 $m$ plus haut. Celui de celle de Chartres est 26 $m$ plus bas que celle de Strasbourg. Quelle est la hauteur de chacun de ces monuments ?

\subsection*{Exercice 3 (4 points)}
Un TGV est composé de 8 wagons de 64 places chacun et d'un wagon-bar de 24 places. À son départ, 29 places sont libres. Combien y a-t-il de passagers? 

\subsection*{Exercice 4 (6 points)}
 
 Au marché, le kilogramme de pommes coûtent 1 euros 80 et celui de tomates coûte 5 euro 40. On achète 700 grammes de pommes et 300 grammes de tomates. 
 \begin{enumerate}
 \item Combien a -t-on payé au total ? 
 \item Quelle denrée a coûté le plus cher ?
 \item On veut désormais acheter au total 1 kg de pommes et tomates de telle manière que l'on ait payé autant pour les pommes que pour les tomates. Quelle masse de pommes et quelle masse de tomates faut-il alors acheter ?
 \end{enumerate}
 	\end{document}
